\section{Kaggle lab: compare supervised baselines fairly}
\label{sec:kaggle-titanic-supervised}

\begin{kaggleproject}{Titanic --- linear versus nonlinear classification}
\kagglesource{https://www.kaggle.com/datasets/yasserh/titanic-dataset}{Titanic Dataset}
\textbf{File.} \texttt{all\_datasets/titanic\_dataset/Titanic-Dataset.csv}\\
\textbf{Target.} \texttt{Survived}.\\
\textbf{Question.} With the same features, development split, preprocessing, and metrics, how do logistic regression and a random forest compare?\\
\textbf{Before running.} Predict which model is more flexible and why a single validation split is not enough to declare a universal winner.

\begin{labmode}
Stage 3A is still workflow-focused. Use the reference pipeline, but you are responsible for adding and interpreting the no-skill baseline.
\end{labmode}

\lstinputlisting[style=mlpython]{all_experiments/lab05-titanic_supervised_baseline.py}

\begin{tryit}
Add a \texttt{DummyClassifier(strategy="most\_frequent")} as a true no-skill baseline. Then explain what additional value logistic regression and the random forest provide over that reference.
\end{tryit}
\end{kaggleproject}

